\chapter{Conclusions and Future Work}

\section{Conclusions}

This thesis demonstrates that dense Channel Knowledge Maps for UAV enabled
networks can be predicted with varying UAV transmitter height, provided that the
problem is treated as a structured propagation task rather than as generic
image to image translation. The final system predicts channel attenuation, delay spread,
and angular spread on \(513\times513\) city maps under a strict city holdout
protocol, with the UAV located at the map centre and its height changing from
sample to sample. The work therefore answers the original practical question
under the CKM simulation contract: a single surrogate can produce dense CKM
targets for unseen cities while
conditioning on continuous UAV transmitter height.

The central lesson is that the strongest model is not a pure black box neural
network. It is a hybrid system in which calibrated physical and statistical priors
explain the stable part of the channel, and a neural residual model corrects
what remains. This division of labour is what made the final result reliable:
the priors capture radial propagation, visibility, height effects, building layout,
and spread scale statistics; the learned model then focuses on local residual
structure and distribution mismatch.

The finished pipeline meets the original numerical targets. The standalone
path loss prior achieved \SI{1.9246}{dB} RMSE on the held out
validation+test split, and \SI{1.9383}{dB} when reused as the frozen prior on
the final 14-city unseen test split. The spread priors reached \SI{28.10}{ns}
delay spread RMSE and \SI{13.74}{\degree} angular spread RMSE on the same final
test split. \textsc{HARP-Net CKM} improved those
anchors to \SI{1.6519}{dB}, \SI{26.5570}{ns}, and
\SI{11.3854}{\degree}, respectively. These numbers are meaningful because the
baseline being improved is already strong and train calibrated, not a weak
baseline.

\subsection{What the final result means}

The clearest contribution is a working CKM predictor for variable height UAV
maps. The final model receives geometry, masks, frozen priors, and a continuous
height condition, and produces dense channel attenuation, delay spread, and
angular spread maps jointly. The height dependence is physically important: the same
city footprint changes its visibility structure, two ray interference pattern,
and spread regimes as the transmitter altitude changes. The final architecture handles
this by combining height aware priors with a sinusoidal height embedding
injected through FiLM conditioning.

The channel attenuation result is the strongest success case. Most LoS path loss
structure is radial, height dependent, and coherent enough that an enhanced
two ray prior can explain much of it before any neural correction is applied.
The residual model still matters, but its role is narrower and more realistic:
it improves local blocked streets, transition regions, and geometry to radio
misalignments that the analytic prior cannot fully capture. Reducing the test
RMSE from \SI{1.9383}{dB} to \SI{1.6519}{dB} after such a strong prior is
therefore a substantial result.

Delay spread and angular spread are harder to judge by a single RMSE. Their
maps contain sparse high value regions, spike like behaviour, and long tailed
outliers. Qualitatively, the final predictions usually recover the main
spatial structure: low spread and high spread regions appear in the right
general parts of the city, and the LoS/NLoS geometry is respected. The
remaining mistakes are mostly local outlier failures. Some groups of pixels
with unusually high DS or AS are missed, while some false positive high spread
patches are predicted where the target remains moderate. This explains why the
spread result is less uniform than the channel attenuation result, while still being useful:
the model improves the global structure and many local regimes, but rare spread
hot spots remain the hardest cases. These RMSE gains should not be read as
universal typical error gains: delay spread MAE increases, and angular spread
MAE is nearly unchanged, so the spread claim remains an RMSE and structure
claim rather than a calibrated error claim.

The most important methodological conclusion is that CKM prediction improves
when the model head follows the target distribution. Earlier point regression
models could produce plausible-looking maps while missing NLoS attenuation modes
or suppressing spread tails. The distribution first stage fixed that diagnosis
by using GMM style outputs and histogram aware evaluation. The final model keeps
that lesson in its head design: it emits mixture parameters around calibrated
anchors and exports their local weighted point prediction, while the final
reported claim remains the deterministic RMSE improvement selected after the
RMSE/MAE dominant fine tuning stage.

\subsection{Answers to the research questions}

\textbf{RQ1. Can a single deep learning surrogate predict per pixel channel attenuation,
delay spread, and angular spread on \(513\times513\) channel knowledge maps,
under strict city holdout evaluation, with errors compatible with the thesis
targets?} Yes, within the fixed CKM contract used in this thesis. The final
model predicts PL, DS, and AS jointly on unseen cities and clears the original
targets: \SI{1.6519}{dB} channel attenuation RMSE, \SI{26.5570}{ns} delay spread RMSE,
and \SI{11.3854}{\degree} angular spread RMSE. The result should be understood
under one carrier frequency, one centred UAV, valid ground receivers only, and
city level train/validation/test separation.

\textbf{RQ2. Which part of the prediction should be left to a calibrated
analytic prior and which part should be left to the neural network, so that the
hybrid model is both accurate and interpretable?} The calibrated priors should
handle the most predictable propagation structure: two ray LoS path loss,
height-binned LoS calibration, morphology aware NLoS corrections, and
log domain ridge models for DS and AS. These priors are strong enough to be
useful on their own. The neural model is then valuable because it learns bounded
residual corrections around them, reducing the frozen-prior test errors from
\SI{1.9383}{dB}, \SI{28.1023}{ns}, and \(13.7416^\circ\) to
\SI{1.6519}{dB}, \SI{26.5570}{ns}, and \(11.3854^\circ\).

\textbf{RQ3. How does the continuous UAV transmitter height
(\SIrange{12}{478}{m}) affect the prediction error, and is the same
architecture able to cover the full height range without per height
retraining?} Height affects the difficulty of the prediction through visibility,
blocking and spread tail behaviour, so it must be represented explicitly. In
the final model, height enters twice: through height aware calibrated priors and
through a sinusoidal FiLM conditioning path inside the neural residual model.
Low altitude cases contain more blocked links, stronger transition zones and
harder spread tails, while high altitude samples are easier even though they
are less represented in the training data. This indicates that the height trend
is physical, not only a sampling artifact.

\subsection{Main contributions}

The first contribution is the final variable height CKM predictor itself:
\textsc{HARP-Net CKM}, a single final CKM model that jointly
predicts channel attenuation, delay spread, and angular spread for unseen city
layouts while conditioning on continuous UAV transmitter height.

The second contribution is the calibrated prior family. The path loss prior
combines free space/radial structure, coherent two ray LoS behaviour,
height dependent calibration, NLoS morphology features, and fallback regimes.
The spread priors extend the same philosophy to DS and AS through log domain
ridge calibration, topology/visibility grouping, height regimes, and robust
fallbacks. Together, these priors put a large amount of physical and empirical
knowledge into a reproducible estimator that is already competitive before the
neural residual is trained.

The third contribution is the distribution first modelling lesson. The thesis
shows that target histograms, LoS/NLoS marginals, and spike plus tail spread
behaviour are not diagnostic extras; they determine which output head is
appropriate. This is why GMM style residual prediction and histogram aware
checks became part of the final design.

The fourth contribution is the evaluation contract: city holdout splitting,
valid ground receiver masking, LoS/NLoS reporting, topology diagnostics, and
transmitter height stratified reporting. This contract is essential because random
splits or unmasked building pixels can make the same model appear stronger
without improving the deployment-relevant task.

\subsection{Practical implications and scope}

The practical implication is that future CKM systems should be compared under
the same contract before architectural claims are made. A new backbone is useful
only if it improves held out city metrics after applying the same receiver mask,
height conditioning, LoS/NLoS reporting, topology analysis, and prior baselines.
Otherwise, a lower error may simply come from solving an easier problem.

The result is strongest as a reproducible surrogate for the CKM simulation
contract used here. It is not a universal substitute for ray tracing reference
generation across all frequencies, antenna patterns, or cities. It does show,
however, that once the dataset contract is fixed, repeated simulator calls can
be compressed into a fast and interpretable predictor. In the Barcelona runtime benchmark, pixel level MATLAB
ray tracing took \SI{99.89}{s} per sample, while the complete final generator
path took \SI{0.88}{s}, a \(113.7\times\) reduction.

\subsection{Negative results}

The negative results are useful mainly as design evidence, so the detailed
attempt history is kept in Appendix~A rather than repeated in the main
conclusion. The short lesson is that pure image translation, larger CNN
backbones, topology experts, multi scale outputs, and uncertainty heads each
exposed a different failure mode, but none removed the need for strong frozen
priors and distribution aware residual learning.

This is why the thesis ends with \textsc{HARP-Net CKM}: the final design is
not just a larger neural network, but a bounded residual predictor built around
the parts of the channel that can already be explained physically or
statistically.

\section{Future Directions}

\paragraph{Short-medium term.}
The next step is to strengthen the current residual pipeline without changing
the evaluation contract, then extend it toward frequency and radio configuration
conditioning. Useful additions include transmitter-oriented mixup, explicit
corridor losses, and controlled ablations of each residual head and loss
weight. The current dataset fixes the carrier and antenna assumptions, while a
more general outdoor CKM model should learn how attenuation and spreads vary with
frequency, bandwidth, antenna pattern, and possibly polarization. A natural
dataset for this stage would simulate the same cities across several radio
configurations so that the model can separate geometry effects from
configuration effects.

A nearer term training direction is adaptation to new or under represented
topology types. The per topology analysis suggests that dense high-rise,
compact mid-rise, open low-rise, mixed urban layouts, and future morphology
classes do not fail in exactly the same way. A practical continuation would
therefore be to fine tune \textsc{HARP-Net CKM} on newly introduced
urban typologies, or to train lightweight specialist heads on top of the shared
backbone when only limited examples are available. This would test whether the
model can recover local structure in topology regimes beyond the original
training distribution, rather than only specialising to already defined classes.

For delay and angular spreads, another short to medium term direction is to
optimize not only pixel-wise RMSE but also spatial structure. Spread maps can
look physically coherent even when small shifts across quantized target
boundaries are penalized harshly by RMSE. Future fine tuning should therefore
test structural objectives, for example gradient consistency, region wise
correlation, SSIM like terms, LoS/NLoS boundary consistency, or topology aware
ranking losses, so that the model preserves the correct spread patterns rather
than only minimizing the average numeric error.

\paragraph{Medium to long term.}
The final extension is field calibration. The attenuation and spread priors are
interpretable, which makes them suitable for simulator to reality correction
using a limited number of real measurements. In the medium to long term, this could
lead to a CKM predictor that starts from ray-traced or prior generated maps,
adapts to a new city with sparse measurements, and still reports uncertainty by
LoS/NLoS regime, transmitter height bin, and topology class. Those safeguards should
remain part of the method before considering operational use in emergency
coverage, UAV network planning, or rapid deployment studies.
